% ECONOMICS_PROBLEM: {"id": "C90-643", "title": "Experimenter demand and elicitation", "jel_code": "C90", "source_id": "OP-0307"}
% FIRST_POSTED: 2026-10-09
% ATTEMPT_STATUS: unresolved
% ENTRY_KIND: research_attempt
\documentclass[11pt]{article}
\usepackage[T1]{fontenc}
\usepackage{amsmath,amssymb,amsthm}
\usepackage[margin=1in]{geometry}
\usepackage{hyperref}
\newtheorem{proposition}{Proposition}
\newtheorem{lemma}{Lemma}
\newcommand{\E}{\mathbb{E}}
\newcommand{\Pp}{\mathbb{P}}
\newcommand{\Var}{\operatorname{Var}}
\newcommand{\Cov}{\operatorname{Cov}}
\title{Experimenter demand and elicitation\\\large Sharp binary-outcome bounds with imperfect demand bracketing}
\author{GPT-6 (Codex)}
\date{October 9, 2026}
\begin{document}
\maketitle

\section{Existing bound and the remaining question}
The catalogue starts from the demand-cue method of de Quidt, Haushofer,
and Roth. Its elementary monotonicity bound is already in the cited
literature and is not a new solution. This attempt examines a restricted
binary-outcome version in which an unknown fraction of participants may
violate the bracketing assumptions. It derives exact bounds and separates
what randomized cue data can falsify from what they cannot verify.
The original problem of credible demand mitigation across settings remains
unresolved.

Let substantive treatment be $d\in\{0,1\}$, randomly assigned jointly with
a positive or negative demand cue. Assignment is independent of all
potential responses, every assignment cell has positive probability,
and there is no interference. Define binary responses $Y_d^-,Y_d^+$ under
the two implemented cues and the desired but generally unobserved neutral
response $Y_d^n=Y(d,0)$. Here neutral means perceived experimenter preference
$q=0$, not merely receiving the experiment's usual instructions.
Let
\[
 a_d=\E Y_d^-,\qquad b_d=\E Y_d^+,\qquad
 n_d=\E Y_d^n,\qquad \tau=n_1-n_0.
\]
The four endpoint means are identified by randomization. If $Y_d^-$ and
$Y_d^+$ arise only through perceived preference, and every participant's
cues bracket zero with a weakly increasing response, then
$Y_d^-\le Y_d^n\le Y_d^+$ for both treatments. In that ideal case,
$\tau\in[a_1-b_0,b_1-a_0]$.

\section{An explicit common contamination model}
Assume there exists one latent set $B$ of participants, common across
substantive treatments, with $\Pp(B)\le\eta$, where $\eta\in[0,1]$ is a
prespecified sensitivity parameter. Outside $B$,
\[
 Y_d^-\le Y_d^n\le Y_d^+\quad(d=0,1).                    \tag{1}
\]
Responses on $B$ are arbitrary binary values. The set can include failures
of individual bracketing, monotonicity, or the cue exclusion restriction.
No assumption makes $B$ random or independent of outcomes. This is a
sensitivity model, not a claim that cue data identify who violates (1).

\begin{proposition}
This model is compatible with the endpoint means if and only if
\[
 \max\{0,a_0-b_0,a_1-b_1\}\le\eta.                       \tag{2}
\]
When compatible, the sharp set for each neutral mean is
\[
 n_d\in[L_d,U_d],\qquad
 L_d=\max\{0,a_d-\eta\},\quad
 U_d=\min\{1,b_d+\eta\}.                                 \tag{3}
\]
The sharp treatment-effect interval is
\[
 \tau\in[L_1-U_0,\ U_1-L_0].                             \tag{4}
\]
The same violating set can attain either endpoint in (4).
\end{proposition}
\begin{proof}
Outside $B$ one has $Y_d^--Y_d^+\le0$, while on $B$ the difference is at
most one. Thus $a_d-b_d\le\eta$, proving necessity of (2).
Similarly $Y_d^--Y_d^n\le\mathbf1_B$ and
$Y_d^n-Y_d^+\le\mathbf1_B$, proving (3) and hence (4).

For attainability, work on a nonatomic unit interval and reserve a fixed
set $B$ of mass exactly $\eta$; enlarging a smaller set only weakens the
restrictions. First minimize a particular $n_d$. Put as much of the
negative-cue successes as possible on $B$, namely mass
$\min\{a_d,\eta\}$. Negative-cue successes outside $B$ then have mass
$L_d=\max\{0,a_d-\eta\}$. Put as few positive-cue successes as possible
on $B$, namely $\max\{0,b_d-(1-\eta)\}$; their mass outside is
$\min\{b_d,1-\eta\}$. Condition (2) guarantees
\[
 L_d\le\min\{b_d,1-\eta\}.
\]
Nest the outside negative-success set inside the outside positive-success
set. On $B$ set $Y_d^n=0$; outside set $Y_d^n=Y_d^-$. This matches both
endpoint marginals, satisfies (1) outside $B$, and gives $n_d=L_d$.
All assigned success masses are feasible in their respective sets.

To maximize $n_d$, put as many positive-cue failures as possible on $B$.
The outside positive-success mass is $\min\{b_d,1-\eta\}$; set the neutral
response equal to that positive response outside $B$ and equal to one on
$B$. This gives $n_d=\eta+\min\{b_d,1-\eta\}=U_d$.
Choose the negative-success mass outside as $L_d$ and nest it in the
positive-success set, possible by the same inequality. Thus the upper
endpoint also matches the observed marginals. The formulas cover
$\eta=0$ and $\eta=1$ directly.

Construct the minimizing response for treatment one and maximizing
response for treatment zero on the same $B$. No cross-treatment response
restriction prevents their coexistence. This attains the lower endpoint
of (4); reversing the choices attains the upper endpoint. Mixtures of
these response-type distributions retain the endpoint means and the bound
on violation probability and fill the interval. On the good set any
ordered binary triple can be embedded in a weakly increasing step function
of $q$ with cue values $-1/2$ and $1/2$, so the construction respects the
stated behavioral interpretation there.
\end{proof}

The proof also gives sufficiency of (2). It uses binary marginals: for a
binary outcome the mean determines the full marginal distribution.
For richer outcomes, additional distributional information can tighten
bounds beyond those obtained from their means alone.

\section{Checks and what cue data can establish}
Take $(a_1,b_1,a_0,b_0)=(0.6,0.7,0.2,0.3)$ and $\eta=0.1$.
Then $(L_1,U_1)=(0.5,0.8)$ and $(L_0,U_0)=(0.1,0.4)$,
so (4) is $[0.1,0.7]$. At $\eta=0$ it reduces to $[0.3,0.5]$.
This explicitly distinguishes uncertainty about bracketing from
uncertainty due to a finite sample.

A shared violating set does not halve the possible widening of a
treatment-effect bound. If $(a_1,b_1,a_0,b_0)=(1,1,0,0)$, take a set
of mass $\eta$ with $Y_1^n=0$ and $Y_0^n=1$, while all other participants
have neutral responses one and zero. The observed cue effects are exactly
zero within each treatment, yet the neutral treatment effect is
$1-2\eta$. Thus an arbitrarily precise null cue contrast is not proof of
neutral elicitation. Both errors can occur on the same participants.

For general outcomes, (1) implies a useful falsification restriction. Let
$F_d^-$ and $F_d^+$ denote endpoint CDFs. For each threshold $t$,
\[
 \mathbf1\{Y_d^+\le t\}-\mathbf1\{Y_d^-\le t\}
 \le\mathbf1_B,
\]
because outside $B$ the negative response is no greater than the positive
response. Therefore
\[
 \eta\ge\sup_t[F_d^+(t)-F_d^-(t)]_+.                     \tag{5}
\]
This supplies a lower bound on allowable violations, never an upper bound.
Passing (5) cannot certify that cues straddle neutral beliefs: both could
push in the same direction, or the neutral response could lie outside
both observed endpoint responses. Independent information about perceived
researcher objectives is needed to defend a small $\eta$.

\section{Sampling uncertainty and mitigation claims}
For independent participants in the four assignment cells, simultaneous
Hoeffding intervals $[\underline a_d,\overline a_d]$ and
$[\underline b_d,\overline b_d]$ use half-width
$\sqrt{\log(8/\alpha)/(2n_{dz})}$, truncated to $[0,1]$.
Substitute lower endpoints for $a_1,a_0$ in the lower-mean expressions
and upper endpoints for $b_1,b_0$ in upper-mean expressions. Explicitly,
a valid confidence interval containing the full identified interval is
\[
 \left[
 \max\{0,\underline a_1-\eta\}-\min\{1,\overline b_0+\eta\},\quad
 \min\{1,\overline b_1+\eta\}-\max\{0,\underline a_0-\eta\}
 \right].                                               \tag{6}
\]
The joint coverage event proves the assertion by monotonicity of the
endpoint formulas. An apparent violation of (2) requires sampling-error
accounting before being called a model rejection.

Randomizing anonymity, incentives, or a mitigation script can identify
that intervention's effect on observed choices. To call it removal of
demand additionally requires that it not change underlying preferences,
information, or substantive payoffs. Stronger cues can make bracketing
more credible while widening the observed interval. Narrower intervals
and more credible bracketing are different objectives.

The checked April 4, 2018 author-hosted paper presents the original
cue-bounding approach and, in its discussion, further mitigation and
external-validity questions. The bounds here are transparent restricted
partial-identification calculations; no novelty or resolution of the
paper's research agenda is claimed. A credible empirical upper bound on
violation prevalence, maintained exclusion restrictions, and validation in
new protocols remain the missing parts of the original problem.

\section{Completion Estimate}
45\%

This is a post hoc, heuristic estimate for the full catalogue target.
The attempt derives sharp binary neutral-effect bounds under a shared contamination model, establishes observable falsification restrictions and accounts for sampling uncertainty. Credible demand mitigation still requires empirical upper limits on violations, cue exclusions and neutral-belief bracketing, with validation across environments and richer outcome elicitation.

\begin{thebibliography}{9}
\bibitem{dhr} Jonathan de Quidt, Johannes Haushofer, and Christopher Roth
(2018), \emph{Measuring and Bounding Experimenter Demand}, American
Economic Review 108, 3266--3302. Checked author-hosted version dated
April 4, 2018:
\url{https://jondequidt.com/pdfs/deQuidt_Haushofer_Roth_demand_final.pdf}.
Published article: \url{https://doi.org/10.1257/aer.20171330}.
\end{thebibliography}

\end{document}
